% m2sv poster -- CTB@ICML 2026
% Portrait, 24in (W) x 36in (H) = 60.96cm x 91.44cm. Build: latexmk -pdf poster.tex
\documentclass[final]{beamer}
\usepackage[orientation=portrait,size=custom,width=60.96,height=91.44,scale=1.05]{beamerposter}
\usepackage{graphicx}
\usepackage{booktabs}
\usepackage{amsmath,amssymb}
\usepackage{tcolorbox}
\tcbuselibrary{skins}

% ---- palette ----
\definecolor{uhgreen}{RGB}{0,76,55}      % deep green
\definecolor{accent}{RGB}{198,146,20}    % gold
\definecolor{ink}{RGB}{30,30,30}
\definecolor{soft}{RGB}{244,246,245}
\definecolor{failbg}{RGB}{251,246,240}
\definecolor{failrule}{RGB}{176,58,46}

\setbeamercolor{block title}{fg=white,bg=uhgreen}
\setbeamercolor{block body}{fg=ink,bg=soft}
\setbeamertemplate{blocks}[rounded][shadow=false]
\setbeamertemplate{navigation symbols}{}
\setbeamerfont{block title}{size=\large,series=\bfseries}

\newcommand{\FIG}{../2026-01-28/figures}

\setlength{\columnsep}{2.0cm}
\newlength{\colw}
\setlength{\colw}{0.475\paperwidth}

% Failure-mode card: #1 map img, #2 sv img, #3 GT/Pred, #4 trace excerpt, #5 reason
\newcommand{\failcard}[5]{%
\begin{tcolorbox}[colback=failbg,colframe=failrule,boxrule=1.2pt,arc=4pt,
  left=10pt,right=10pt,top=8pt,bottom=8pt,title=#3,
  coltitle=white,colbacktitle=failrule,fonttitle=\normalsize\bfseries]
\begin{minipage}[c]{0.30\linewidth}\centering
  \includegraphics[width=\linewidth]{#1}
\end{minipage}\hfill
\begin{minipage}[c]{0.30\linewidth}\centering
  \includegraphics[width=\linewidth]{#2}
\end{minipage}\hfill
\begin{minipage}[c]{0.36\linewidth}\footnotesize
  \textit{``#4''}
\end{minipage}
\vspace{4pt}\par\footnotesize\textbf{Why:} #5
\end{tcolorbox}\vspace{0.5cm}}

\begin{document}
\begin{frame}[t]

% ================= TITLE BAR =================
\begin{tcolorbox}[colback=uhgreen,colframe=uhgreen,boxrule=0pt,arc=8pt,left=18pt,right=18pt,top=14pt,bottom=14pt]
\centering
{\color{white}\Huge\textbf{m2sv: A Scalable Benchmark for}}\\[6pt]
{\color{white}\Huge\textbf{Map-to-Street-View Spatial Reasoning}}\\[12pt]
{\color{white}\large Yosub Shin \quad Michael Buriek \quad Igor Molybog}\\[4pt]
{\color{white}\normalsize University of Hawai\textquoteright i at M\=anoa\ \ (Buriek also PwC)}\\[6pt]
{\color{accent}\large\textbf{CTB Workshop @ ICML 2026 \quad$\cdot$\quad non-archival}}\\[4pt]
{\color{white}\normalsize \texttt{github.com/YosubShin/m2sv} \qquad live demo: \texttt{m2sv.yosubshin.com}}
\end{tcolorbox}

\vspace{0.9cm}

% ================= COLUMNS =================
\begin{columns}[t]

% ===================== LEFT COLUMN: what & how well =====================
\begin{column}{\colw}

\begin{block}{The Task}
Given a \textbf{north-up overhead map} with labeled candidate directions and a
\textbf{Street View photo} taken at the same intersection, infer \emph{which
labeled direction the camera was facing}. The camera sits at the intersection
center; letters point outward along each road.
\vspace{0.5cm}
\begin{center}
\includegraphics[width=0.44\linewidth]{\FIG/sample-question-1-map.jpg}\hfill
\includegraphics[width=0.44\linewidth]{\FIG/sample-question-1-sv.jpg}
\end{center}
\vspace{0.2cm}
\footnotesize\emph{Solve by aligning stable map geometry (roads, parcels,
landmarks) with the ground-level view --- not transient cues.}
\end{block}

\begin{block}{Benchmark \& Data (open)}
\begin{itemize}
  \item \textbf{m2sv-20k} --- 20{,}000 real intersections (10k train / 10k
        validation) across \textbf{32 cities}, multiple continents, controlled
        candidate ambiguity.
  \item \textbf{m2sv-sft-11k} --- 11k curated structured reasoning traces for SFT.
  \item Released as \textbf{blueprints} (coordinates + azimuths) plus rendering
        code to respect map licensing; all eval/training code public.
\end{itemize}
\end{block}

\begin{block}{Zero-shot: humans clear the best VLM}
\centering
\begin{tabular}{l r r}
\toprule
Model & N & Acc.\ (\%)\\
\midrule
Gemini-3-Pro            & 1k  & 65.2 $\pm$ 3.0\\
GPT-5                   & 1k  & 57.2 $\pm$ 3.1\\
Gemini-2.5-Pro          & 1k  & 47.2 $\pm$ 3.1\\
Qwen3-VL-235B (Think)   & 1k  & 42.7 $\pm$ 3.1\\
Random baseline         & 1k  & 31.4 $\pm$ 2.9\\
\midrule
\textbf{Human (engaged, $n{=}8$)} & 200 & \textbf{72.0 $\pm$ 9.7}\\
\textbf{Human (expert)}            & 200 & \textbf{95.0}\\
\bottomrule
\end{tabular}
\vspace{0.4cm}
\flushleft\footnotesize Multi-annotator study (12 participated; the 10 who finished
all 200 are analyzed). Two low-agreement outliers excluded by \emph{agreement}
($\kappa_{\text{expert}}\!<\!0.3$), not accuracy.
\end{block}

\begin{block}{Human baseline \& agreement}
\begin{center}
\includegraphics[width=0.92\linewidth]{\FIG/human_accuracy_dist.png}
\end{center}
\footnotesize Attentive humans converge (median pairwise $\kappa=0.48$, up to
$0.76$); Gemini-3-Pro sits \emph{below} the human consensus (median $\kappa=0.42$),
Qwen3-VL-235B lower still ($0.15$). The gap shows up in \emph{agreement structure},
not only accuracy.
\end{block}

\begin{block}{Difficulty: the hardest-tier collapse}
\begin{center}
\includegraphics[width=0.86\linewidth]{\FIG/human_difficulty_decircular.png}
\end{center}
\footnotesize Difficulty bins from \textbf{human response time}; accuracy from a
model-independent axis. Humans degrade gently ($83\!\to\!63\%$) and stay above both
models; models lose the most ground on the hardest tier, falling to (Qwen) or
toward (Gemini, $50\%$) chance. \textbf{Symmetric} intersections (azimuth-gap ratio
$S=\frac{\max}{\min}$) are hardest for all models.
\end{block}

\end{column}

% ===================== RIGHT COLUMN: why models fail =====================
\begin{column}{\colw}

\begin{block}{Why models fail (qualitative analysis)}
\footnotesize Humans solve nearly all of these. Model errors come from
\emph{brittle alignment strategies and reliance on unstable evidence}, not missing
information. Each card: overhead map, street view, a trace excerpt, and the error.
\end{block}

\vspace{0.4cm}

\failcard{\FIG/failure-mode-0-map.jpg}{\FIG/failure-mode-0-sv.jpg}
{Ego$\leftrightarrow$allo inversion \hfill GT: A\ \ Pred: C}
{Looking from the center towards A, the road curves to the right --- contradicting the left curve in the photo.}
{The map shows a \emph{left} curve, but the model reasons as if it were a right curve. Left/right inversion at the map$\leftrightarrow$photo boundary.}

\failcard{\FIG/failure-mode-1-map.jpg}{\FIG/failure-mode-1-sv.jpg}
{Over-reliance on unstable cues \hfill GT: B\ \ Pred: D}
{Both corners appear to have gray roofs; there is no reddish house on the left corner, so B is incorrect.}
{Roof color is lighting- and angle-dependent and weakly represented in overhead imagery. A human would discount this cue.}

\failcard{\FIG/failure-mode-5-map.jpg}{\FIG/failure-mode-5-sv.jpg}
{Landmark misbinding \hfill GT: A\ \ Pred: B}
{The corner house has a large swimming pool in its backyard; this does not match the street view, so A is wrong.}
{The corner house has \emph{no} pool --- one exists two houses down. The model binds a nearby landmark to the wrong spatial anchor.}

\failcard{\FIG/failure-mode-2-map.jpg}{\FIG/failure-mode-2-sv.jpg}
{Symmetry trap \hfill GT: D\ \ Pred: B}
{Several headings look equally plausible; the model commits on a weak cue instead of resolving the ambiguity.}
{At symmetric intersections, models fail to \emph{escalate} to fine cues (lane counts, markings, curb geometry) that would disambiguate.}

\begin{block}{Takeaways}
\begin{itemize}
  \item m2sv isolates a scalable \textbf{map$\leftrightarrow$street-view alignment}
        primitive frontier VLMs have \emph{not} solved (best 65\% vs human 72/95\%).
  \item Humans beat the best model \emph{and} agree with one another; models lie
        outside that consensus and collapse on the hardest, symmetric cases.
  \item SFT$+$RL on m2sv helps ($+$10 pts) but transfer is limited --- failures are
        about \emph{adapting reasoning to ambiguity}, not missing data.
\end{itemize}
\vspace{0.3cm}
\centering\small\textbf{Code/data:} \texttt{github.com/YosubShin/m2sv} \qquad
\textbf{Try it:} \texttt{m2sv.yosubshin.com}
\end{block}

\end{column}

\end{columns}

\vspace{0.5cm}
\begin{center}
\footnotesize\color{ink} Supported by NSF NRT-AI 2244574. \quad Contact: \texttt{yosubs@hawaii.edu}
\end{center}

\end{frame}
\end{document}
